\documentclass[11pt]{article}
\usepackage[margin=0.9in]{geometry}
\usepackage{amsmath,amssymb,amsthm,booktabs,mathrsfs}
\usepackage[T1]{fontenc}
\usepackage{lmodern,microtype}
\usepackage[colorlinks=true,linkcolor=blue,urlcolor=blue]{hyperref}
\newtheorem{theorem}{Theorem}[section]
\newtheorem{lemma}[theorem]{Lemma}
\newtheorem{proposition}[theorem]{Proposition}
\newtheorem{corollary}[theorem]{Corollary}
\theoremstyle{remark}
\newtheorem{remark}[theorem]{Remark}
\newcommand{\norm}[1]{\left\|#1\right\|}
\title{Toward a four-dimensional Yang--Mills construction\\
\large Notebook 33: gauge-invariant determinant densities and Cartan probes}
\author{Research working document}
\date{17 October 2026}
\begin{document}
\maketitle
\begin{abstract}
We identify a constant-curvature calculation with the actual
Maxwell coefficient of the leading quantum correction.
The starting point is the exact vector, based-ghost and
effective-Hessian determinant identity of Notebook 32.

At a fixed composite endpoint, the matrices transform by
background-gauge conjugation. Their diagonal color traces
therefore define individually gauge-invariant analytic
densities. Fixed-endpoint gaps and exponentially localized
source kernels make their matrix logarithms local, uniformly
in compatible volume. A rooted-density Ward argument proves
that their quadratic affine coefficient is precisely the
second-displacement moment of the total Hessian. This
resolves the ambiguity that would arise from evaluating
an arbitrary integration-by-parts density on an unbounded
affine connection.

Commuting affine connections are exactly Wilson stationary
and are preserved by the actual centered ordinary averaging
maps. An exponential-kernel uniqueness argument identifies
their polynomial stationary jets with the constructed
minimum section. The envelope identity extends the
zero-current conclusion to the effective classical actions.
The constraint multiplier consequently vanishes through
the background orders needed for the Maxwell calculation.

The resulting coefficient is an explicit diagonal trace
expression on a constant Cartan-curvature probe, using
the actual composite constraint and its normalized
determinants. For the bare composite it involves the
Wilson Hessian directly. The coarse contribution already
has a uniform bound. We compute a specified massive
continuum proper-time reference, including its positive
logarithmic coefficient and a uniform finite-part bound.
An exact heat-integral formula for the actual lattice
coefficient isolates sufficient estimates on three
heat-time intervals. Those estimates for the actual
family remain incomplete. A closed-walk expansion proves
the short-time bound for the actual based-ghost part,
uniformly in endpoint and preceding depth. The vector
part, the other heat-time intervals and the subsequent
full-measure and continuum arguments remain.
\end{abstract}

All new arguments below are working proofs checked against
the stated inputs from Notebooks 15--32. The matrix-logarithm
locality estimates fix the endpoint before taking a volume
limit. They are not estimates uniform in the ultraviolet
depth. An affine field is used only through polynomial
extensions of localized derivatives; no periodic family
of continuously varying magnetic flux is assumed.
